\textbf{Phần 1. Trắc nghiệm lựa chọn}
\setcounterpageref{ex}{0}

\begin{ex} %[5D6V34]
Một khu rừng hình chữ nhật có chiều dài $2{,}5\text{ km}$ và chiều rộng $1{,}6\text{ km}$. Diện tích của khu rừng đó tính theo đơn vị héc-ta là
\choice
{$4\text{ ha}$}
{$40\text{ ha}$}
{\True $400\text{ ha}$}
{$4\,000\text{ ha}$}
\loigiai{
Diện tích khu rừng theo ki-lô-mét vuông là:
$2{,}5 \times 1{,}6 = 4\text{ (km}^2\text{)}$.\\
Vì $1\text{ km}^2 = 100\text{ ha}$ nên $4\text{ km}^2 = 400\text{ ha}$.\\
\textbf{Đáp án của câu hỏi là C}
}
\end{ex}

\begin{ex} %[5D6V34]
Một xe ô tô chở $3{,}5$ tấn gạo và $18$ tạ ngô. Tổng khối lượng hàng hóa xe ô tô đó đang chở tính theo đơn vị ki-lô-gam là
\choice
{$3\,680\text{ kg}$}
{$4\,300\text{ kg}$}
{\True $5\,300\text{ kg}$}
{$5\,800\text{ kg}$}
\loigiai{
Đổi sang ki-lô-gam:
$3{,}5\text{ tấn} = 3\,500\text{ kg}$; $18\text{ tạ} = 1\,800\text{ kg}$.\\
Tổng khối lượng hàng hóa là:
$3\,500 + 1\,800 = 5\,300\text{ (kg)}$.\\
\textbf{Đáp án của câu hỏi là C}
}
\end{ex}

\begin{ex} %[5D6V34]
Điền dấu thích hợp vào chỗ chấm: $4\text{ m}^2\text{ }5\text{ dm}^2 \dots 4{,}5\text{ m}^2$.
\choice
{\True $<$}
{$>$}
{$=$}
{Không so sánh được}
\loigiai{
Đổi sang cùng đơn vị mét vuông:
$4\text{ m}^2\text{ }5\text{ dm}^2 = 4{,}05\text{ m}^2$.\\
Vì $4{,}05 < 4{,}50$ nên $4\text{ m}^2\text{ }5\text{ dm}^2 < 4{,}5\text{ m}^2$.\\
\textbf{Đáp án của câu hỏi là A}
}
\end{ex}

\begin{ex} %[5D6V34]
Một người đi xe đạp từ A đến B dài $18\text{ km}$. Sau khi đi được $12{,}5\text{ km}$ thì người đó dừng lại nghỉ. Quãng đường còn lại phải đi tính theo đơn vị mét là
\choice
{$550\text{ m}$}
{\True $5\,500\text{ m}$}
{$6\,500\text{ m}$}
{$55\,000\text{ m}$}
\loigiai{
Quãng đường còn lại là:
$18 - 12{,}5 = 5{,}5\text{ (km)}$.\\
Đổi sang mét: $5{,}5\text{ km} = 5\,500\text{ m}$.\\
\textbf{Đáp án của câu hỏi là B}
}
\end{ex}

\begin{ex} %[5D6V34]
Một cánh đồng hình chữ nhật có chiều dài $800\text{ m}$ và chiều rộng $500\text{ m}$. Cứ mỗi héc-ta thu hoạch được $6$ tấn thóc. Toàn bộ cánh đồng đó thu hoạch được bao nhiêu tấn thóc?
\choice
{$120\text{ tấn}$}
{$180\text{ tấn}$}
{\True $240\text{ tấn}$}
{$300\text{ tấn}$}
\loigiai{
Diện tích cánh đồng là:
$800 \times 500 = 400\,000\text{ (m}^2\text{)}$.\\
Đổi: $400\,000\text{ m}^2 = 40\text{ ha}$.\\
Số tấn thóc thu hoạch được là:
$40 \times 6 = 240\text{ (tấn)}$.\\
\textbf{Đáp án của câu hỏi là C}
}
\end{ex}

\begin{ex} %[5D6V34]
Một bể chứa nước có dung tích $3{,}5\text{ m}^3$. Hiện tại trong bể đang có $1\,800\text{ l}$ nước. Cần bơm thêm vào bể bao nhiêu lít nước nữa để đầy bể?
\choice
{$1\,500\text{ l}$}
{\True $1\,700\text{ l}$}
{$1\,800\text{ l}$}
{$2\,000\text{ l}$}
\loigiai{
Đổi: $3{,}5\text{ m}^3 = 3\,500\text{ dm}^3 = 3\,500\text{ l}$.\\
Số lít nước cần bơm thêm là:
$3\,500 - 1\,800 = 1\,700\text{ (l)}$.\\
\textbf{Đáp án của câu hỏi là B}
}
\end{ex}

\begin{ex} %[5D6V34]
Một tấm bìa có diện tích $0{,}25\text{ m}^2$. Người ta cắt lấy một phần có diện tích $650\text{ cm}^2$. Diện tích phần bìa còn lại là
\choice
{$1\,750\text{ cm}^2$}
{\True $1\,850\text{ cm}^2$}
{$1\,950\text{ cm}^2$}
{$2\,050\text{ cm}^2$}
\loigiai{
Đổi: $0{,}25\text{ m}^2 = 2\,500\text{ cm}^2$.\\
Diện tích phần bìa còn lại là:
$2\,500 - 650 = 1\,850\text{ (cm}^2\text{)}$.\\
\textbf{Đáp án của câu hỏi là B}
}
\end{ex}

\begin{ex} %[5D6V34]
Có ba bao đường: bao thứ nhất nặng $45{,}5\text{ kg}$, bao thứ hai nặng hơn bao thứ nhất $3{,}8\text{ kg}$, bao thứ ba nhẹ hơn bao thứ hai $5{,}3\text{ kg}$. Tổng khối lượng cả ba bao đường là
\choice
{$135{,}6\text{ kg}$}
{\True $138{,}8\text{ kg}$}
{$140{,}2\text{ kg}$}
{$142{,}5\text{ kg}$}
\loigiai{
Khối lượng bao thứ hai là:
$45{,}5 + 3{,}8 = 49{,}3\text{ (kg)}$.\\
Khối lượng bao thứ ba là:
$49{,}3 - 5{,}3 = 44\text{ (kg)}$.\\
Tổng khối lượng cả ba bao là:
$45{,}5 + 49{,}3 + 44 = 138{,}8\text{ (kg)}$.\\
\textbf{Đáp án của câu hỏi là B}
}
\end{ex}

\begin{ex} %[5D6V34]
Một mảnh đất hình vuông có chu vi đo được là $1\,200\text{ m}$. Diện tích của mảnh đất đó bằng bao nhiêu héc-ta?
\choice
{$6\text{ ha}$}
{\True $9\text{ ha}$}
{$12\text{ ha}$}
{$90\text{ ha}$}
\loigiai{
Độ dài cạnh mảnh đất hình vuông là:
$1\,200 : 4 = 300\text{ (m)}$.\\
Diện tích mảnh đất là:
$300 \times 300 = 90\,000\text{ (m}^2\text{)}$.\\
Đổi: $90\,000\text{ m}^2 = 9\text{ ha}$.\\
\textbf{Đáp án của câu hỏi là B}
}
\end{ex}

\begin{ex} %[5D6V34]
Phép chuyển đổi đơn vị đo lường nào sau đây là SAI?
\choice
{$2\text{ km }45\text{ m} = 2{,}045\text{ km}$}
{$5\text{ tấn }8\text{ kg} = 5{,}008\text{ tấn}$}
{\True $3\text{ m}^2\text{ }4\text{ dm}^2 = 3{,}4\text{ m}^2$}
{$7\text{ ha }50\text{ m}^2 = 7{,}005\text{ ha}$}
\loigiai{
Vì $1\text{ m}^2 = 100\text{ dm}^2$ nên $4\text{ dm}^2 = 0{,}04\text{ m}^2$.\\
Do đó $3\text{ m}^2\text{ }4\text{ dm}^2 = 3{,}04\text{ m}^2 
eq 3{,}4\text{ m}^2$.\\
\textbf{Đáp án của câu hỏi là C}
}
\end{ex}

\begin{ex} %[5D6V34]
Một thợ may dùng cuộn chỉ dài $150\text{ m}$. Biết may một chiếc áo sơ mi hết trung bình $12{,}5\text{ m}$ chỉ. Cuộn chỉ đó may được nhiều nhất số chiếc áo sơ mi là
\choice
{$10\text{ chiếc}$}
{\True $12\text{ chiếc}$}
{$14\text{ chiếc}$}
{$15\text{ chiếc}$}
\loigiai{
Số chiếc áo sơ mi may được là:
$150 : 12{,}5 = 12\text{ (chiếc)}$.\\
\textbf{Đáp án của câu hỏi là B}
}
\end{ex}

\begin{ex} %[5D6V34]
Một vận động viên chạy 3 vòng quanh một hồ nước hình tròn có bán kính $100\text{ m}$. Quãng đường người đó đã chạy là bao nhiêu mét?
\choice
{$628\text{ m}$}
{$1\,256\text{ m}$}
{\True $1\,884\text{ m}$}
{$3\,140\text{ m}$}
\loigiai{
Chu vi một vòng quanh hồ là:
$2 \times 3{,}14 \times 100 = 628\text{ (m)}$.\\
Quãng đường chạy 3 vòng là:
$628 \times 3 = 1\,884\text{ (m)}$.\\
\textbf{Đáp án của câu hỏi là C}
}
\end{ex}

\textbf{Phần 2. Trắc nghiệm đúng sai}
\setcounterpageref{ex}{0}

\begin{ex} %[5D6V34]
Xét tính đúng sai của các phép chuyển đổi đơn vị đo lường sau:
\choiceTF
{\True $15\text{ tấn }8\text{ tạ} = 15{,}8\text{ tấn}$}
{\True $6\text{ km}^2\text{ }45\text{ ha} = 6{,}45\text{ km}^2$}
{Số đo $2\text{ m}^2\text{ }7\text{ cm}^2 = 2{,}07\text{ m}^2$}
{\True Độ dài $3{,}5\text{ m} = 350\text{ cm}$}
\loigiai{
\begin{itemchoice}
\item Đúng. $8\text{ tạ} = 0{,}8\text{ tấn}$ nên $15\text{ tấn }8\text{ tạ} = 15{,}8\text{ tấn}$.
\item Đúng. Vì $1\text{ km}^2 = 100\text{ ha}$ nên $45\text{ ha} = 0{,}45\text{ km}^2$.
\item Sai. Vì $1\text{ m}^2 = 10\,000\text{ cm}^2$ nên $7\text{ cm}^2 = 0{,}0007\text{ m}^2$, vậy $2\text{ m}^2\text{ }7\text{ cm}^2 = 2{,}0007\text{ m}^2 
eq 2{,}07\text{ m}^2$.
\item Đúng. Vì $1\text{ m} = 100\text{ cm}$ nên $3{,}5\text{ m} = 350\text{ cm}$.
\end{itemchoice}
}
\end{ex}

\begin{ex} %[5D6V34]
Một đoàn xe chở hàng gồm 3 xe tải lớn và 2 xe tải nhỏ. Mỗi xe tải lớn chở $4{,}5$ tấn hàng, mỗi xe tải nhỏ chở $2{,}5$ tấn hàng.
\choiceTF
{\True Tổng khối lượng hàng do 3 xe tải lớn chở là $13{,}5$ tấn}
{\True Tổng khối lượng hàng do 2 xe tải nhỏ chở là $5$ tấn}
{\True Cả đoàn xe chở được tất cả $18{,}5$ tấn hàng, tương đương $185$ tạ}
{Khối lượng hàng trung bình mỗi xe chở trong đoàn là $4$ tấn}
\loigiai{
\begin{itemchoice}
\item Đúng. Khối lượng hàng xe lớn chở: $3 \times 4{,}5 = 13{,}5\text{ (tấn)}$.
\item Đúng. Khối lượng hàng xe nhỏ chở: $2 \times 2{,}5 = 5\text{ (tấn)}$.
\item Đúng. Tổng khối lượng: $13{,}5 + 5 = 18{,}5\text{ (tấn)} = 185\text{ tạ}$.
\item Sai. Trung bình mỗi xe chở: $18{,}5 : 5 = 3{,}7\text{ (tấn)} 
eq 4\text{ tấn}$.
\end{itemchoice}
}
\end{ex}

\textbf{Phần 3. Trắc nghiệm trả lời ngắn}
\setcounterpageref{ex}{0}

\begin{ex} %[5D6V34]
Một khu công nghiệp hình chữ nhật có chiều dài $2\text{ km}$ và chiều rộng $1{,}5\text{ km}$. Diện tích của khu công nghiệp đó là bao nhiêu héc-ta?
\shortans{300}
\loigiai{
Diện tích khu công nghiệp là:
$2 \times 1{,}5 = 3\text{ (km}^2\text{)}$.\\
Đổi: $3\text{ km}^2 = 300\text{ ha}$.
}
\end{ex}

\begin{ex} %[5D6V34]
Một bao gạo nặng $50\text{ kg}$. Một xe tải chở $80$ bao gạo như thế. Tổng khối lượng gạo mà xe tải đó đang chở là bao nhiêu tấn?
\shortans{4}
\loigiai{
Tổng khối lượng gạo xe tải chở là:
$50 \times 80 = 4\,000\text{ (kg)}$.\\
Đổi: $4\,000\text{ kg} = 4\text{ tấn}$.
}
\end{ex}

\begin{ex} %[5D6V34]
Một mảnh vườn hình chữ nhật có chiều dài $40\text{ m}$ và chiều rộng $25\text{ m}$. Cứ mỗi $100\text{ m}^2$ người ta thu hoạch được $80\text{ kg}$ rau. Hỏi cả mảnh vườn thu hoạch được bao nhiêu yến rau?
\shortans{80}
\loigiai{
Diện tích mảnh vườn là:
$40 \times 25 = 1\,000\text{ (m}^2\text{)}$.\\
Khối lượng rau thu hoạch được là:
$1\,000 : 100 \times 80 = 800\text{ (kg)}$.\\
Đổi: $800\text{ kg} = 80\text{ yến}$.
}
\end{ex}

\begin{ex} %[5D6V34]
Một cuộn dây cáp quang dài $1\text{ km}$. Người ta cắt thành các đoạn nhỏ bằng nhau, mỗi đoạn dài $25\text{ m}$. Hỏi cắt được tất cả bao nhiêu đoạn dây cáp như thế?
\shortans{40}
\loigiai{
Đổi: $1\text{ km} = 1\,000\text{ m}$.\\
Số đoạn dây cáp cắt được là:
$1\,000 : 25 = 40\text{ (đoạn)}$.
}
\end{ex}

\begin{ex} %[5D6V34]
Một khu rừng hình vuông có diện tích bằng $16\text{ ha}$. Chu vi của khu rừng đó bằng bao nhiêu mét?
\shortans{1600}
\loigiai{
Đổi: $16\text{ ha} = 160\,000\text{ m}^2$.\\
Vì $400 \times 400 = 160\,000$ nên độ dài cạnh khu rừng là $400\text{ m}$.\\
Chu vi của khu rừng là:
$400 \times 4 = 1\,600\text{ (m)}$.
}
\end{ex}

\begin{ex} %[5D6V34]
Một chiếc xe ô tô có bình xăng chứa đầy được $60\text{ l}$. Sau một chuyến đi, xe đã tiêu thụ hết $20\text{ l}$ xăng. Trong bình xăng của xe còn lại bao nhiêu lít xăng?
\shortans{40}
\loigiai{
Lượng xăng còn lại trong bình là:
$60 - 20 = 40\text{ (l)}$.
}
\end{ex}
